%!TEX root = main.tex
%=========================================================

\section{Evaluation}
\label{sec:evaluation}

\er{
Goal for this section: show experimentally that the proposed algorithms and protocols answer the research questions set forth in the introduction.
We aim at answering the following research questions:
\begin{itemize}
	\item Does the use of a single protocol in an heterogeneously distributed MANET lead to sub-optimal performance?
	\item Is it beneficial to use different protocols for broadcast in a heterogeneously-distributed MANET, and does it reduce contention/collisions, costs, and improves overall reliability?
	\item What is the relative impact of heterogeneity in motion, and heterogeneity in placement, to the difference in performance and reliability observed in the complete MANET?
	\item What is the cost of the interoperability mechanism when a subset of the network decides to adapt and emerge an overlay in its region?
	\item Are the observables available locally to each node a good representation of the actual deployment environment of each node?
	\item How effective are the adaptation strategies in autonomously deciding on the emergence of overlays or the use of flooding-based broadcast?
\end{itemize}
The section can be structured as follows:
\begin{itemize}
	\item Experimental setup: simulator, traces, models for the nodes energy consumption, etc.
	\item Recap that we use our MAC layer
	\item Microbenchmarks with the square setup
	\item Macrobenchmarks with more complex setups, with different dense/dynamic regions over space
\end{itemize}
}
